\documentclass[11pt]{article}
\usepackage[utf8]{inputenc}
\usepackage[T1]{fontenc}
\usepackage{geometry}
\usepackage{hyperref}
\geometry{margin=1in}


\title{Trabajo de Laboratorio 02: Pruebas sobre el comportamiento de la memoria caché}
\author{Cesar Lengua Malaga}
\date{\today}

\begin{document}
\maketitle

\section*{Resumen}
Se implementaron: (1) los dos pares de bucles del libro para multiplicacion matriz-vector, (2) multiplicacion de matrices clasica (tres bucles) y (3) multiplicacion por bloques (seis bucles) en C++ y Go. Se comparo el tiempo de ejecucion, el comportamiento de cache y la complejidad. Se usan Cachegrind y KCachegrind para medir cache misses.

\section*{Metodologia}
Se escribieron programas simples en C++ y Go para ejecutar los tres casos. Los tiempos se midieron con cronometro interno. Para cache se uso Valgrind/Cachegrind y se visualizo con KCachegrind.

\section*{Implementacion}
\noindent 1) Dos bucles anidados (cap. 2, pag. 22). \\
El libro propone multiplicacion matriz-vector con dos ordenes de bucles:
\begin{verbatim}
/* First pair of loops */
for (i = 0; i < MAX; i++)
  for (j = 0; j < MAX; j++)
    y[i] += A[i][j] * x[j];

/* Second pair of loops */
for (j = 0; j < MAX; j++)
  for (i = 0; i < MAX; i++)
    y[i] += A[i][j] * x[j];
\end{verbatim}
En C++/Go se implementaron como \texttt{orden i-j} y \texttt{orden j-i}.
En el programa de matrices se uso un tamaño de bloque fijo de 32 para la version por bloques.

\noindent 2) Multiplicacion clasica. \\
Version de tres bucles anidados:
\begin{verbatim}
for i
  for j
    for k
      C[i][j] += A[i][k] * B[k][j]
\end{verbatim}

\noindent 3) Multiplicacion por bloques. \\
Se agregan tres bucles externos para trabajar por sub-bloques. Esto mejora la localidad de datos al reutilizar sub-bloques en cache.

\section*{Movimiento de datos y cache}
Las matrices estan en memoria en orden por filas. Por eso, acceder en forma contigua (por filas) suele ser mas rapido que acceder por columnas, porque se aprovecha mejor la cache.
Matriz-vector (orden i-j vs orden j-i):
\begin{itemize}
  \item orden i-j: A se recorre por filas (contiguo). Buen uso de cache.
  \item orden j-i: A se recorre por columnas (saltos de n). Peor localidad.
  \item El vector \texttt{x[j]} se reutiliza bien en ambos casos.
\end{itemize}

Clasica (i-j-k):
\begin{itemize}
  \item A se recorre por filas (contiguo en memoria). Buen uso de cache.
  \item B se recorre por columnas (saltos de \texttt{n}). Esto genera muchos cache misses.
  \item C se escribe por filas, pero cada \texttt{C[i][j]} se vuelve a tocar en cada k.
\end{itemize}

Bloques:
\begin{itemize}
  \item Se toman sub-bloques A(i0,k0) y B(k0,j0) y se reutilizan varias veces.
  \item Si el bloque cabe en cache, se reducen accesos a memoria principal.
  \item El algoritmo sigue siendo O(n\textsuperscript{3}) pero con menos fallos de cache.
\end{itemize}

\section*{Complejidad}
El costo asintotico explica el crecimiento del tiempo cuando aumenta N. En la practica, la cache cambia el rendimiento real, por eso se comparan los ordenes de bucle y la version por bloques.
\begin{itemize}
  \item Matriz-vector: O(n\textsuperscript{2}).
  \item Multiplicacion clasica: O(n\textsuperscript{3}).
  \item Multiplicacion por bloques: O(n\textsuperscript{3}) con mejor localidad.
\end{itemize}

\section*{Comandos de compilacion y ejecucion}
C++
\begin{verbatim}
# Compilar
 g++ -O2 -std=c++17 cpp/2loops.cpp -o cpp/2loops_cpp
 g++ -O2 -std=c++17 cpp/matmul.cpp -o cpp/matmul_cpp

# Matriz-vector (orden i-j y orden j-i)
 ./cpp/2loops_cpp

# Multiplicacion de matrices (clasica y por bloques)
 ./cpp/matmul_cpp
\end{verbatim}

Go
\begin{verbatim}
# Compilar
 go build -o go/2loops_go go/2loops.go
 go build -o go/matmul_go go/matmul.go

# Matriz-vector (orden i-j y orden j-i)
 ./go/2loops_go

# Multiplicacion de matrices (clasica y por bloques)
 ./go/matmul_go
\end{verbatim}

\section*{Ejecucion paso a paso}
Matriz-vector (dos bucles, paso a paso):
\begin{verbatim}
# C++ (puedes cambiar el N)
./cpp/2loops_cpp 256
./cpp/2loops_cpp 512
./cpp/2loops_cpp 1024

# Go
./go/2loops_go 256
./go/2loops_go 512
./go/2loops_go 1024
\end{verbatim}

Matrices (clasica y bloques, paso a paso):
\begin{verbatim}
# C++ (N por lista)
./cpp/matmul_cpp 64 128 256 384

# Go (N por lista)
./go/matmul_go 64 128 256 384
\end{verbatim}
La salida que interesa son los misses de cache (I1, D1 y LL). 

\section*{Cachegrind y KCachegrind}
\begin{verbatim}
# C++ con Cachegrind
 valgrind --tool=cachegrind --cachegrind-out-file=cachegrind.out.%p \
   ./cpp/matmul_cpp 256

# Ver resumen en terminal
 cg_annotate cachegrind.out.<pid> | head -n 40

# Abrir en KCachegrind (GUI)
 kcachegrind cachegrind.out.<pid>
\end{verbatim}
Cachegrind simula caches L1/L2 y reporta misses por linea de codigo.
KCachegrind permite explorar esos datos de forma visual.

\section*{Cache misses}
Medicion con Cachegrind para N=256 y N=384 (C++).
\begin{center}
\begin{tabular}{|c|c|c|c|c|c|}
\hline
N & Bloque & I1mr & D1mr & ILmr & DLmr \\
\hline
256 & 16 & 2,477 & 27,778,382 & 2,302 & 9,348 \\
256 & 32 & 2,232 & 33,923,364 & 2,121 & 9,308 \\
256 & 64 & 2,477 & 33,823,110 & 2,302 & 9,348 \\
384 & 16 & 2,478 & 58,963,118 & 2,303 & 9,348 \\
384 & 32 & 2,236 & 77,578,628 & 2,124 & 9,308 \\
384 & 64 & 2,457 & 114,550,754 & 2,284 & 9,348 \\
\hline
\end{tabular}
\end{center}

\section*{Resultados}
C++ (ms) - Matriz-vector:
\begin{center}
\begin{tabular}{|c|c|c|c|}
\hline
Caso & n & Tiempo \\
\hline
Matriz-vector (orden i-j) & 256 & 0.053 \\
Matriz-vector (orden j-i) & 256 & 0.079 \\
Matriz-vector (orden i-j) & 512 & 1.189 \\
Matriz-vector (orden j-i) & 512 & 1.955 \\
Matriz-vector (orden i-j) & 1024 & 1.093 \\
Matriz-vector (orden j-i) & 1024 & 3.302 \\
\hline
\end{tabular}
\end{center}

C++ (ms) - Matrices con distintos bloques: \\
Bloque = 16
\begin{center}
\begin{tabular}{|c|c|c|c|}
\hline
N & Clásica (ms) & Bloques (ms) & Razón \\
\hline
64 & 3.051 & 3.506 & 0.87 \\
128 & 28.619 & 30.382 & 0.94 \\
256 & 262.217 & 263.410 & 1.00 \\
384 & 972.126 & 830.145 & 1.17 \\
512 & 2647.530 & 1910.379 & 1.39 \\
768 & 8467.531 & 6977.349 & 1.21 \\
\hline
\end{tabular}
\end{center}

Bloque = 32
\begin{center}
\begin{tabular}{|c|c|c|c|}
\hline
N & Clásica (ms) & Bloques (ms) & Razón \\
\hline
64 & 3.904 & 3.891 & 1.00 \\
128 & 25.506 & 27.741 & 0.92 \\
256 & 240.560 & 212.774 & 1.13 \\
384 & 773.093 & 703.198 & 1.10 \\
512 & 2161.254 & 1727.578 & 1.25 \\
768 & 6967.306 & 5533.900 & 1.26 \\
\hline
\end{tabular}
\end{center}

Bloque = 64
\begin{center}
\begin{tabular}{|c|c|c|c|}
\hline
N & Clásica (ms) & Bloques (ms) & Razón \\
\hline
64 & 11.289 & 9.897 & 1.14 \\
128 & 34.904 & 23.689 & 1.47 \\
256 & 215.911 & 192.825 & 1.12 \\
384 & 766.689 & 747.056 & 1.03 \\
512 & 1936.274 & 1627.846 & 1.19 \\
768 & 5887.355 & 5273.868 & 1.12 \\
\hline
\end{tabular}
\end{center}

Go (ms) - Matriz-vector:
\begin{center}
\begin{tabular}{|c|c|c|c|}
\hline
Caso & n & Tiempo \\
\hline
Matriz-vector (orden i-j) & 256 & 0.162 \\
Matriz-vector (orden j-i) & 256 & 0.098 \\
Matriz-vector (orden i-j) & 512 & 0.579 \\
Matriz-vector (orden j-i) & 512 & 0.506 \\
Matriz-vector (orden i-j) & 1024 & 2.577 \\
Matriz-vector (orden j-i) & 1024 & 7.665 \\
\hline
\end{tabular}
\end{center}

Go (ms) - Matrices con distintos bloques: \\
Bloque = 16
\begin{center}
\begin{tabular}{|c|c|c|c|}
\hline
N & Clásica (ms) & Bloques (ms) & Razón \\
\hline
64 & 0.636 & 0.708 & 0.90 \\
128 & 7.382 & 9.915 & 0.74 \\
256 & 44.309 & 42.065 & 1.05 \\
384 & 107.165 & 143.984 & 0.74 \\
512 & 264.598 & 351.651 & 0.75 \\
768 & 961.025 & 1096.208 & 0.88 \\
\hline
\end{tabular}
\end{center}

Bloque = 32
\begin{center}
\begin{tabular}{|c|c|c|c|}
\hline
N & Clásica (ms) & Bloques (ms) & Razón \\
\hline
64 & 0.280 & 0.521 & 0.54 \\
128 & 2.247 & 4.265 & 0.53 \\
256 & 33.968 & 39.235 & 0.87 \\
384 & 105.833 & 131.557 & 0.80 \\
512 & 281.786 & 307.529 & 0.92 \\
768 & 1164.062 & 1103.962 & 1.05 \\
\hline
\end{tabular}
\end{center}

Bloque = 64
\begin{center}
\begin{tabular}{|c|c|c|c|}
\hline
N & Clásica (ms) & Bloques (ms) & Razón \\
\hline
64 & 0.519 & 0.647 & 0.80 \\
128 & 7.535 & 5.764 & 1.31 \\
256 & 45.618 & 36.890 & 1.24 \\
384 & 94.441 & 113.721 & 0.83 \\
512 & 277.589 & 300.832 & 0.92 \\
768 & 1069.625 & 951.485 & 1.12 \\
\hline
\end{tabular}
\end{center}

\section*{Analisis de resultados}
En C++, el orden i-j fue consistentemente mas rapido que j-i en matriz-vector, lo que coincide con el acceso contiguo por filas. En Go se observa que para N=256 y N=512 el orden j-i fue ligeramente mas rapido, pero para N=1024 el orden i-j vuelve a ser mejor. Esto puede deberse a efectos del runtime, la cache y el sistema operativo.

Para matrices, C++ muestra mejora con bloques cuando N aumenta, y el bloque 32 o 64 suele dar mejores tiempos que 16. En Go el resultado depende del bloque: con bloque 16 casi siempre gana la clasica, con bloque 32 hay casos mixtos y con bloque 64 se observa mejora en algunos tamaños (por ejemplo N=128, 256 y 768).

La columna Razón se define como tiempo clásica dividido entre tiempo bloques. Si la razón es mayor que 1, la versión por bloques es más rápida; si es menor que 1, la versión clásica es más rápida.

\section*{Conclusiones}
El orden de los bucles cambia la localidad y los cache misses. En matriz-vector, el orden i-j suele ser mas rapido que el orden j-i. En matrices, la version por bloques puede ser mas rapida cuando el tamaño del bloque es adecuado y N es grande, como se ve en los tiempos de C++ para bloques 32 y 64. En Go no siempre mejora, lo cual muestra que el rendimiento depende del lenguaje, del bloque y del tamaño. En ambos lenguajes se mantiene O(n\textsuperscript{3}), pero la cache explica parte de las diferencias observadas.

\section*{Codigo}

Los codigos se encuentran en el siguiente link: \url{https://github.com/rxsh/Parallel} 

\end{document}
